% =====================================================================
%  Phase C --- The station catalogue and the generated world
%  SCOPE: one source of coordinates; the world is built from it.
% =====================================================================
\section{The Station Catalogue and the Generated World}
\label{sec:phaseC_catalogue}

\subsection{Forty-eight stations}

The \phaseB{} greenhouse (Sec.~\ref{sec:phaseB_twin}) is
$9.90 \times 4.90$~m internally, with three planting gutters 0.40~m wide
centred at $y = -1.20$, $0.00$ and $+1.20$~m, and eight plants per row
spaced 0.90~m apart from $x = -3.15$~m. \phaseC{} adds a
\emph{measurement station} on each side of each plant: 3 rows
$\times$ 8 plants $\times$ 2 sides $= 48$, labelled \station{Pi,jR} and
\station{Pi,jL}.

A station is a point on the floor where the robot's sensor head must
stand, offset \SI{0.55}{m} across the row from the plant centre. The
label grammar is deliberately the one an agronomist would speak:
\station{P2,5} names a \emph{plant} and expands to both of its stations;
\station{P2,5R} names one side. Expansion happens on the Cloud, before
the order is signed, so the robot never interprets a wildcard --- a
robot that expanded \texttt{ALL} itself would silently disagree with the
Cloud the day the greenhouse gained a row.

\subsection{The decision that shapes everything else}

Three programs need these coordinates and must not disagree: the world
generator that paints the marks, the robot that drives onto them, and
the dashboard that draws them.

The obvious construction is to type the numbers into the world file and
again into the robot. That is how a mark ends up painted at $y=-1.70$
and driven to at $y=-1.65$, with the robot then reporting that it is
``on'' a mark it is 5~cm away from and nothing anywhere to catch it.
\phaseC{} therefore inverts the dependency: the numbers exist once, in
\filename{agri/catalogue.py}, and \textbf{the greenhouse is generated
from them}.

\begin{lstlisting}[style=bash]
python3 -m agri.world.make_world   # reads the Phase B scene,
                                   # paints 48 marks from the catalogue
python3 -m agri.world.make_robot   # adds the boom and floor camera
\end{lstlisting}

The regression suite closes the loop in both directions: it reads the
plant positions back out of the \phaseB{} world and asserts the
catalogue still agrees with them, and it parses the generated SDF and
asserts every painted mark is where the catalogue says. Neither the
greenhouse nor the robot can drift from the map without a check going
red (Sec.~\ref{sec:phaseC_regression}).

The generator also refuses to run twice over its own output, because
re-painting an already-painted world produces 96 marks and a scene that
looks subtly wrong rather than obviously broken.

\subsection{Two geometric constraints that fixed the numbers}

Two constants were not chosen by preference. They are reported here
because both were got wrong first, and because each is the kind of value
that looks arbitrary until the reason is written down.

\paragraph{Station reach, \SI{0.55}{m}} An inner aisle is only 0.80~m
wide between two gutter edges. The natural design --- turn the base to
face the plant --- points the chassis' long axis across that aisle,
putting 0.29~m of half-length toward the gutter instead of 0.19~m of
half-width, and the two stations sharing an inner aisle then cannot both
fit with usable margin. The first attempt left \SI{1}{cm}.

\phaseC{} keeps the base \emph{along} the aisle and lets the existing
pan head look sideways, which is what the head is for. That converts
0.29 into 0.19 and buys back 0.20~m. At a reach of \SI{0.55}{m}, the
tightest of all 48 stations (\station{P2,1R}) then clears every gutter
and wall by \measured{\SI{0.16}{m}}, and the suite asserts that bound on
every run.

\paragraph{Cross arm, \SI{0.04}{m}} The price of the above is that the
two stations sharing an inner aisle sit only \SI{0.10}{m} apart, at the
same $x$. Their $y$-arms therefore lie on one line. Any arm longer than
\SI{0.05}{m} makes the two painted crosses touch, and a blob detector
then sees \emph{one} marker straddling both --- it would centre the
robot on the midpoint, \SI{0.05}{m} from either station, and file the
visit under whichever label it guessed.

\SI{0.04}{m} leaves a \SI{0.02}{m} gap, about ten pixels for the floor
camera. An earlier version used \SI{0.09}{m} and carried a comment
asserting the crosses did not merge; rendering the pair and running the
detector over it is what showed the comment was wrong. The suite now
renders adjacent pairs and asserts the nearer mark wins at three
different offsets.

\subsection{Why the marks are painted at all}

\SI{0.10}{m} is inside the error a drifting odometric estimate
accumulates in a single lap of this greenhouse. Driving to a station on
odometry alone is therefore not merely imprecise --- it is capable of
parking the robot closer to the neighbouring station than to the one it
is about to file a reading under, which corrupts the \emph{label} rather
than the value. The painted marks exist so that the final approach can
be closed visually against a physical feature of the building, which is
the subject of Sec.~\ref{sec:phaseC_align}.
